\subsection{An additional full profinite failure of the pentagon}

\begin{result}{B.2.1}{Proved}{The kernel of the CM character contains elements failing the pentagon}
There is an element $g\in\GTz$ of multiplier one with
$\Delta(g)=0$ which fails the full profinite pentagon.
Its image in the free pro-$11$ group begins in lower-central
degree nine. Together with A.7, this gives the strict inclusion
$\GT\subsetneq\ker\Delta$. The assertion checks and supplies
an independent pentagon certificate for the existence part of
LM7\_PUSH D.52; it does not use that return's unprovided
degree-nine kernel computation.
\end{result}
\begin{proof}
We give the integral vector, the actual lifting argument and a
small exact image of its pentagon. Write $L(x,y)$ for the free
Lie algebra over $\mathbb Z[1/2]$, and use standard Lyndon
bracketing: a Lyndon word $w$ is bracketed as
$[\mathsf L(u),\mathsf L(v)]$, where $v$ is its longest proper
Lyndon suffix and $w=uv$. Let $2h\in L_9(x,y)$ have the following
coefficients; every unlisted coefficient is zero.
\[
\begin{array}{lr|lr|lr}
 w&[w](2h)&w&[w](2h)&w&[w](2h)\\ \hline
 xxxxxxyyy&2&xxxxxyxyy&-18&xxxxxyyxy&-24\\
 xxxxxyyyy&-6&xxxxyxxyy&18&xxxxyxyxy&-12\\
 xxxxyxyyy&19&xxxxyyxxy&18&xxxxyyxyy&48\\
 xxxxyyyxy&49&xxxxyyyyy&6&xxxyxxxyy&-2\\
 xxxyxxyxy&32&xxxyxxyyy&-1&xxxyxyxxy&20\\
 xxxyxyxyy&17&xxxyxyyxy&11&xxxyxyyyy&-13\\
 xxxyyxxyy&-20&xxxyyxyxy&-12&xxxyyxyyy&-43\\
 xxxyyyxxy&-20&xxxyyyxyy&-105&xxxyyyyxy&-49\\
 xxxyyyyyy&-2&xxyxxyxyy&-41&xxyxxyyxy&-72\\
 xxyxxyyyy&-17&xxyxyxxyy&59&xxyxyxyxy&-18\\
 xxyxyxyyy&-34&xxyxyyxyy&-18&xxyxyyyxy&-68\\
 xxyxyyyyy&12&xxyyxxyyy&-11&xxyyxyxyy&-51\\
 xxyyxyyyy&30&xxyyyxyxy&-10&xxyyyxyyy&40\\
 xxyyyyxyy&60&xxyyyyyxy&24&xyxyxyxyy&6
\end{array}                                               \tag{B.W1}
\]
The notation $[w]$ in the table means the coefficient of
$\mathsf L(w)$, rather than an associative-word coefficient.
Put $z=-x-y$ and $k=-h(z,x)$. Expansion in the free associative
algebra proves
\[
\begin{aligned}
 h(x,y)+h(y,x)&=0,\\
 [x,h(x,y)]&=[x+y,h(z,y)],\\
 [h(x,y),y]&=[k,x+y],\\
 h(x,y)+h(y,z)+h(z,x)&=0.
\end{aligned}                                             \tag{B.W2}
\]
The program \texttt{check\_weight9\_witness.py} computes the
standard brackets recursively, expands all $412$ nonzero
associative coefficients of $2h$, and checks every coefficient
of all four remainders. The saved JSON gives the full source,
the companion $2k$ and the zero remainders. It also computes
the right-$y$ Fox abelianization as
$2X^6Y^2+6X^5Y^3+6X^4Y^4+2X^3Y^5$.
Consequently the actual quotient by the second derived Lie ideal
sends $h$ to $X^2Y^2(X+Y)^3[x,y]$.

We next justify the pro-$11$ realization, which was outside the
parts of D.44 checked by the supplied referee. Let $P$ be the
free pro-$11$ group on $x,y$ and put $z=(xy)^{-1}$ in the group.
The integral lower-central graded Lie algebra is the free
$\mathbb Z_{11}$ Lie algebra. This follows directly from the
Hall collection model: the free nilpotent quotient has the
unique ordered basic-commutator coordinates, its pro-$11$
completion has the same coordinates over $\mathbb Z_{11}$,
and the commutator's first term is the integral Hall bracket.
Passing over all nilpotent classes recovers $P$, since every
finite $11$-group kills a sufficiently deep lower-central term.

The derivation $D_h(x)=0$, $D_h(y)=[h,y]$ is tangential and
special by the third identity of (B.W2). It has an integral
special lift. To prove this, start with group conjugators at
$y,z$ realizing the prescribed degree-nine terms, keeping $x$
fixed. The product error in $xy'z'$ starts above their
prescribed special relation. At degree $m$, modifying the
two conjugators in degree $m-1$ changes that error by
\[
          (u,v)\longmapsto[u,y]+[v,-x-y].                 \tag{B.W3}
\]
Its image is $[L_{m-1},x]+[L_{m-1},y]=L_m$.
The last equality is integral: Jacobi moves one generator
to the outer position in each bracket monomial, by induction
on the length of its first bracket factor. We can therefore
cancel each successive error, preserving every lower term.
The conjugators converge in every finite $11$-quotient.
The resulting continuous endomorphism fixes all three
peripheral conjugacy classes and acts identically on the
Frattini quotient. It is an automorphism.

Let $I^9$ be the degree-nine subgroup of the IA special
outer automorphism group. It is pro-$11$, as is its closed
filtration kernel: the IA automorphism group of a finitely
generated pro-$11$ group is an inverse limit of finite
$11$-groups, and taking closed subgroups and compact
quotients preserves this property. Construction (B.W3)
makes its leading map onto the integral special outer
derivations. The actual geometric $\mathfrak S_3$ acts
on $P$ through outer automorphisms and acts on $I^9$
by conjugation. Equations (B.W2) show, for its two
generators $s(x,y)=(y,x)$ and $r(x,y)=(y,z)$, that
\[
 sD_hs^{-1}-D_h=-\operatorname{ad}(h),\qquad
 rD_hr^{-1}-D_h=-\operatorname{ad}(h).
\]
Thus $-[D_h]$ is fixed in the actual leading outer quotient.

Fixed points lift across this equivariant pro-$11$
surjection. Indeed, lift a fixed target element to $b$.
In the semidirect product with $\mathfrak S_3$, the
two complements obtained from the original copy of
$\mathfrak S_3$ and its conjugate by $b$ lie over the
same finite quotient. Coprime complement conjugacy in
the pro-$11$ kernel supplies a kernel element which
adjusts $b$ to commute with that copy. The finite
Schur--Zassenhaus assertion applies at every finite
kernel quotient; the closed nonempty sets of such
adjustments have the finite intersection property,
so compactness gives one adjustment for all quotients.
We obtain an actual special outer automorphism
$\beta_{11}$ with leading class $-[D_h]$, centralizing
the same geometric $\mathfrak S_3$.

For completeness we retain the precise source normalization.
Put $B_{11}=\widehat B/O^{11}(F)$, where
$F=\ker(\widehat B\to\mathfrak S_3)$, with
$x=\tau^{-2},y=sxs$ and $\tau=sr$.
Its pure kernel is the same $P$. Since $P$ is centerless,
the extension is canonically the pullback
\[
 B_{11}\simeq\operatorname{Aut}(P)
       \mathbin{\times_{\operatorname{Out}(P)}}\mathfrak S_3,
 \quad b\longmapsto(\operatorname{Ad}_b,\bar b).           \tag{B.W4}
\]
Injectivity follows from $Z(P)=1$; a lift of $\bar b$
and an inner adjustment prove surjectivity. Centerlessness
follows from $C_P(x)=\langle x\rangle$ and
$C_P(y)=\langle y\rangle$, whose intersection is trivial.
The primitive-generator centralizer follows in every
nilpotent class by comparing the first nonzero Hall term,
and then by the inverse limit.

Choose a representative $\psi$ of $\beta_{11}$ fixing
$x$ literally. Conjugation by $\psi$ on the first
factor of (B.W4) gives an automorphism $\phi$ over
$\mathfrak S_3$. Since $\phi(\tau)^2=\tau^2=x^{-1}$
and both elements centralize $x$, write
$\phi(\tau)=x^a\tau$; taking squares gives $2a=0$
in $\mathbb Z_{11}$, hence $\phi(\tau)=\tau$.
The images of $s,r$ are conjugate to their original
positive generators by coprime complement conjugacy.
For $f_1=s\phi(s)$ the abelianization is $(a,-a)$,
because $\phi(s)^2=1$ and $s$ interchanges $x,y$.
Adjusting $\phi$ by $\operatorname{Inn}(x^a)$ replaces
$f_1$ with $y^af_1x^{-a}\in P'$. Denote it by $f_{11}$.
The literal torsion equations now give both gentle
cycles, and the restricted pure action is
\[
             x\longmapsto x,\qquad
             y\longmapsto f_{11}^{-1}yf_{11}.             \tag{B.W5}
\]

The normalization retains the prescribed leading term.
In $P/\gamma_{10}P$, centralizing $x$ means belonging
to $\langle x\rangle\gamma_9P/\gamma_{10}P$.
To verify this, remove its abelian $x$-power and
consider its first remaining term of degree below nine.
A homogeneous associative polynomial commuting with $x$
is a multiple of a pure power of $x$; such a term has
zero coefficient in a Lie monomial of degree at least two.
There is therefore no such first term. Apply this to
an inner representative of (B.W5) modulo $\gamma_{10}$.
Its equality on $y$ and the zero abelianization of
$f_{11}$ force both remaining peripheral powers to be
zero, so $f_{11}\in\gamma_9P$. The leading derivation
of (B.W5) is $-D_{h_f}$. Equality of its outer class
with $-[D_h]$, followed by the same free-Lie
centralizer calculation, gives $h_f=h$.

Use the refereed relative lifting theorem LM7\_PUSH D.50
on its actual source
\[
 \widehat B/R_{\rm nil}(F)
   \simeq\prod_{p,\mathfrak S_3}\widehat B/O^p(F).
 \tag{B.W6}
\]
Place $\phi$ in the $11$-component and the identity in
every other component. Their quotient actions on
$\mathfrak S_3$ agree. At $11$ the positive $s,r$
conjugators lie in $P$ by the same coprime argument;
elsewhere they are one, and the cusp conjugator is one
in every component. These verify the exact peripheral
hypotheses of D.50, rather than just its graded target.
It gives an actual full profinite gentle element
$g_0=(1,f)$ with $f\mapsto f_{11}$.

We now test its pentagon through an actual graded
comparison. The lower-central Lie algebra of the pure
four-strand group, tensored with $\mathbb Q_{11}$,
is $\mathfrak t_4$, with generators $t_{ij}$ and
relations
$[t_{ij},t_{kl}]=0$ for disjoint pairs and
$[t_{ij},t_{ik}+t_{jk}]=0$ for distinct indices.
The quotient by the full twist is
$\mathfrak p_5=\mathfrak t_4/\langle\sum_{i<j}t_{ij}\rangle$;
these are the actual lower-central comparison maps in
\cite[Section 5.1.1 and Remark 5.1]{EnriquezFurusho2021}.
One can see compatibility with the pro-$11$ projection
directly on the successive free kernels of pure-strand
forgetting: their integral Hall lattices have torsion-free
graded pieces and complete to those same lattices over
$\mathbb Z_{11}$. The braid insertion maps send the
degree-one classes to the named sums of $t_{ij}$.
As $f_{11}$ begins in degree nine, the first term of
its group pentagon is exactly
\[
\begin{split}
 \mathcal P(h)={}&h(t_{12},t_{23}+t_{24})
                 +h(t_{13}+t_{23},t_{34})\\
 &-h(t_{23},t_{34})
   -h(t_{12}+t_{13},t_{24}+t_{34})
   -h(t_{12},t_{23}).                                    
\end{split}\tag{B.W7}
\]
Products become sums at this first degree, and conjugating
an insertion changes only higher terms.

Here a $24$-dimensional representation suffices. In
$\mathbb Q_{11}[\mathfrak S_4]$ put
$C=\sum_{i<j}(ij)$ and send
$t_{ij}\mapsto(ij)-C/6$. The class sum $C$ is central.
Disjoint transpositions commute, and multiplication in
each three-letter symmetric subgroup proves
$[(ij),(ik)+(jk)]=0$. The sum of all six images is zero.
This defines a Lie homomorphism out of the very sphere
quotient $\mathfrak p_5$ just used. Adding central
elements to the arguments of a homogeneous Lie word
of degree at least two changes none of its values.
Exact multiplication in the integral group algebra
therefore gives
\[
          \mathcal P(2h)\longmapsto
                    29120\bigl((34)-(12)\bigr)\ne0.
 \tag{B.W8}
\]
The program checks every infinitesimal braid relation
and all $24$ coefficients of this remainder.
This is a necessary image of the full group pentagon:
a zero group pentagon would have zero first graded
term and zero image under this homomorphism. Hence
$g_0$ fails the full profinite pentagon.

Finally take the square $g=g_0^2$ in the group $\GTz$.
R1 gives $\Delta(g)=2\Delta(g_0)=0$.
The multiplier-one composition law in (B.W5) adds
first degree-nine terms, so the first term of its
pro-$11$ pair is $2h$; substitutions by the degree-nine
IA action change that term only in higher degree.
Its pentagon still has exactly the nonzero image
(B.W8). Thus $g\in\ker\Delta\setminus\GT$, as asserted.
\end{proof}

\begin{result}{B.2}{Open}{Finite images and the determined quadratic annihilator}
The complete finite-image comparison remains unresolved.
The group
\[
 \mathscr X=
 \{\chi\in\operatorname{Hom}_{\rm cts}(\GTz,\mathbb F_2):
                     \chi|_{\GT}=0\}
\]
has dimension exactly $\aleph_0$, by B.2.24--B.2.25,
and (B.I11) gives an exact basis isomorphism from the
countable direct sum of $\mathbb F_2$. Every character
$\Delta\circ D_P^n$ belongs to this space, and these
characters are linearly independent even on $H$.
We have exhibited the
finite characteristic detector A.4.2, and proved
$\GT\subsetneq\ker\Delta$ in B.2.1.
B.2.2 determines the actual images in the free
exponent-eleven quotients through class ten, with the
first separation at class eight. B.2.3 supplies the
smaller integral weight-seven witness. B.2.4 proves
that the action on all odd-prime free pro-$\ell$
quotients cannot detect a nonzero element of
$\mathscr X$.
B.2.5--B.2.8 compute a full marked central cover and
use the exact finite-certificate lifting theorem to
construct a second independent character
$\Delta\circ D_P$. Thus generation of $\mathscr X$
by $\Delta$ is refuted. B.2.25 additionally computes
the full gentle and GT images in every finite prefix
quotient of the actual binary shift map.  B.2.27 proves
that this particular descent family does not span
$\mathscr X$, by an actual full descent-kernel element
detected by the harmonic character.  B.2.28 gives two
jointly independent character families and computes
the full image of every finite prefix of their paired
shift quotient.  The determination of the abstract
group and dimension retains the exact full-basis map.
\end{result}

The same character group is equivalently the continuous
$\mathbb F_2$ dual of
\[
 \GTz\big/
 \overline{\langle[\GTz,\GTz],\ (\GTz)^2,\ \GT\rangle}.
 \tag{B.W9}
\]
Indeed every quadratic character kills the three displayed
subgroups, and every continuous linear functional on their
closed quotient pulls back to a character with the required
vanishing. B.2.25 proves that this quotient is the
countable binary product, by evaluation on a full
chosen basis. Formula (B.W9) is its actual quotient map.
The tests undertaken here are the signed
CM action, its congruence-cover homology refinement, and
the full weight-nine lift above, the exact low-weight
finite images, and the odd-prime Frattini comparison.
The remaining finite-image task includes the common
five-point census of A.4 and explicit images at further
mixed characteristic quotients. The proved abstract
classification of the quadratic quotient supplies the
maps (B.I3) at every finite prefix, while these further
geometric finite-image comparisons remain to be computed.
